\documentclass[11pt]{article}

\usepackage[T1]{fontenc}
\usepackage[utf8]{inputenc}
\usepackage{lmodern}
\usepackage{microtype}
\usepackage[a4paper,margin=30mm]{geometry}
\usepackage{amsmath,amssymb,amsthm,mathtools}
\usepackage{enumitem}
\usepackage{xcolor}
\usepackage{hyperref}

\hypersetup{
  colorlinks=true,
  linkcolor=blue!55!black,
  urlcolor=blue!65!black,
  citecolor=blue!55!black,
  pdftitle={Digit independence under a common multiplier},
  pdfauthor={Mathematical idea by Mikhail Katkov}
}

\newtheorem{theorem}{Theorem}
\newtheorem{lemma}{Lemma}
\newtheorem{remark}{Remark}

\newcommand{\N}{\mathbb{N}}
\newcommand{\Z}{\mathbb{Z}}
\newcommand{\Zmod}[1]{\Z/#1\Z}
\newcommand{\digit}{d}
\newcommand{\digits}{D}
\newcommand{\LI}{\operatorname{LI}}

\title{Digit Independence under a Common Multiplier}
\author{Mathematical idea and construction: Mikhail Katkov\thanks{Weizmann Institute of Science; Institute for Advanced Studies; \href{mailto:mikhail.katkov@gmail.com}{mikhail.katkov@gmail.com}.}\\
  \small Formalization and exposition developed with the assistance of\\
  \small ChatGPT 5.6 (OpenAI)  \\
  \small LLM orcherstration: Nick Velken Krom}
\date{August 2026}

\begin{document}

\maketitle

\begin{abstract}
For every integer base $p\ge 2$ and every finite set $S$ of distinct positive
natural numbers, we construct a single multiplier $K$ such that the infinite,
zero-padded base-$p$ digit vectors of the numbers $aK$, with $a\in S$, are
linearly independent over $\Zmod p$. The base need not be prime. The proof
constructs a unit lower-triangular matrix from threshold-crossing digits and
then realizes its rows as linear projections of digit vectors produced by one
block-packed multiplier. A Lean formalization accompanies this note.
\end{abstract}

\section{Digit vectors and the theorem}

Fix a natural-number base $p\ge 2$. For $N,r\in\N$, define the $r$th digit of
$N$, viewed in $\Zmod p$, by
\begin{equation}
  \digit_{p,r}(N)
  =
  \left[
    \left\lfloor \frac{N}{p^r}\right\rfloor
  \right]_p
  \in \Zmod p.
  \label{eq:digit}
\end{equation}
Equivalently, this is the ordinary $r$th base-$p$ digit
$\lfloor N/p^r\rfloor\bmod p$, followed by the canonical map into
$\Zmod p$.

The infinite, zero-padded digit vector is
\begin{equation}
  \digits_p(N)
  =
  \bigl(
    \digit_{p,0}(N),\digit_{p,1}(N),\digit_{p,2}(N),\ldots
  \bigr)
  \in (\Zmod p)^{\N}.
  \label{eq:digit-vector}
\end{equation}

\begin{theorem}[Digit independence under a common multiplier]
\label{thm:main}
Let $p\ge 2$, and let $S\subset\N_{>0}$ be finite. There exists $K\in\N$
such that the family
\[
  \bigl(\digits_p(aK)\bigr)_{a\in S}
\]
is linearly independent over $\Zmod p$.
\end{theorem}

Pairwise distinctness is built into the fact that $S$ is a set. No primality
assumption is made. Thus Theorem~\ref{thm:main} includes composite $p$, even
though $\Zmod p$ then has zero divisors.

In Lean the family is indexed by the subtype of elements of a
\texttt{Finset Nat}. The condition that zero does not belong to the finset
expresses positivity. The main theorem is named
\texttt{DigitIndependence.exists\_digits\_linearIndependent}.

\section{Proof}

Write
\begin{equation}
  A=\sum_{a\in S}a,
  \qquad
  m=A^2+1,
  \qquad
  P=p^m.
  \label{eq:parameters}
\end{equation}
Because $p\ge2$, the elementary estimate $x<2^{x+1}$ gives
\begin{equation}
  A^2<P.
  \label{eq:large-P}
\end{equation}

If $S$ is empty, the assertion is immediate. The construction below also
handles that case formally. We therefore proceed with the estimates needed
for elements of $S$.

\subsection{First threshold crossings}

For every $i\in S$, let $q_i$ be the least natural number for which
\begin{equation}
  P\le iq_i.
  \label{eq:crossing}
\end{equation}
Such a number exists because $i>0$. Minimality yields
\begin{equation}
  i(q_i-1)<P\le iq_i.
  \label{eq:minimal-crossing}
\end{equation}
Moreover $q_i\le P$, because $i\ge1$ implies $P\le iP$.

\subsection{A unit lower-triangular matrix}

Define a square matrix indexed by $S$:
\begin{equation}
  B_{a,i}
  =
  \digit_{p,m}(aq_i)
  -
  \digit_{p,m}\bigl(a(q_i-1)\bigr)
  \in\Zmod p.
  \label{eq:B}
\end{equation}

We first compute its diagonal. From
Equation~\eqref{eq:minimal-crossing},
\[
  P\le iq_i=i(q_i-1)+i<P+i<2P.
\]
The final inequality follows from $i\le A<P$. Since $P=p^m$, division by
$p^m$ shows that
\[
  \digit_{p,m}(iq_i)=1,
  \qquad
  \digit_{p,m}\bigl(i(q_i-1)\bigr)=0.
\]
Consequently,
\begin{equation}
  B_{i,i}=1.
  \label{eq:diagonal}
\end{equation}

Now suppose $a<i$. From Equation~\eqref{eq:minimal-crossing} we have
$iq_i<P+i$. Using commutativity of multiplication,
\begin{align*}
  i(aq_i)
  &=a(iq_i)\\
  &<a(P+i)\\
  &<(a+1)P\\
  &\le iP.
\end{align*}
The strict middle inequality is equivalent to $ai<P$, which follows from
$a,i\le A$ and Equation~\eqref{eq:large-P}. The last inequality follows from
$a+1\le i$. Cancelling the positive factor $i$ gives
\begin{equation}
  aq_i<P=p^m.
  \label{eq:off-diagonal-bound}
\end{equation}
The same bound holds for $a(q_i-1)$. Hence both digits in
Equation~\eqref{eq:B} vanish, and
\begin{equation}
  B_{a,i}=0
  \qquad\text{whenever }a<i.
  \label{eq:above-zero}
\end{equation}

After ordering $S$ by the usual order on $\N$,
Equations~\eqref{eq:diagonal} and~\eqref{eq:above-zero} say that $B$ is lower
triangular with every diagonal entry equal to one. Therefore
\begin{equation}
  \det B=1.
  \label{eq:det-one}
\end{equation}

This equality is essential for composite bases. Over a ring with zero
divisors, a merely nonzero determinant need not imply invertibility. Here the
determinant is exactly $1$, hence a unit. It follows that $B$ is invertible and
its rows are linearly independent over $\Zmod p$.

\subsection{Packing all threshold multipliers into one number}

It remains to replace the separate multipliers $q_i$ by one common multiplier
$K$. Choose an enumeration
\[
  S=\{s_0,\ldots,s_{n-1}\}.
\]
Choose a block length $L$ such that
\begin{equation}
  m<L,
  \qquad
  aq_i<p^L
  \quad\text{for all }a,i\in S.
  \label{eq:block-bound}
\end{equation}
The Lean proof uses the explicit, deliberately generous value
\begin{equation}
  L=AP+m+1.
  \label{eq:L}
\end{equation}
Indeed $a\le A$ and $q_i\le P$, so $aq_i\le AP<p^L$.

Pack both $q_i$ and $q_i-1$ into disjoint blocks of $L$ base-$p$ digits:
\begin{equation}
  K
  =
  \sum_{j=0}^{n-1}q_{s_j}p^{jL}
  +
  \sum_{j=0}^{n-1}(q_{s_j}-1)p^{(n+j)L}.
  \label{eq:K}
\end{equation}
Multiplication distributes over this expression:
\[
  aK
  =
  \sum_{j=0}^{n-1}(aq_{s_j})p^{jL}
  +
  \sum_{j=0}^{n-1}a(q_{s_j}-1)p^{(n+j)L}.
\]
Condition~\eqref{eq:block-bound} ensures that every coefficient fits in its
own block. Therefore no carry crosses a block boundary, and
\begin{align}
  \digit_{p,jL+m}(aK)
    &=\digit_{p,m}(aq_{s_j}),
    \label{eq:first-block}\\
  \digit_{p,(n+j)L+m}(aK)
    &=\digit_{p,m}\bigl(a(q_{s_j}-1)\bigr).
    \label{eq:second-block}
\end{align}

Define a linear coordinate projection
\begin{equation}
  T:(\Zmod p)^{\N}\longrightarrow(\Zmod p)^S
\end{equation}
by
\begin{equation}
  (Tx)_{s_j}
  =x_{jL+m}-x_{(n+j)L+m}.
  \label{eq:T}
\end{equation}
Equations~\eqref{eq:B}, \eqref{eq:first-block}, and
\eqref{eq:second-block} give
\begin{equation}
  T\bigl(\digits_p(aK)\bigr)=B_{a,\bullet},
  \label{eq:T-row}
\end{equation}
where $B_{a,\bullet}$ denotes the row of $B$ indexed by $a$.

The rows of $B$ are linearly independent. If the original vectors
$\digits_p(aK)$ satisfied a linear relation, applying $T$ would give the same
relation among the rows of $B$. All its coefficients must therefore vanish.
This proves Theorem~\ref{thm:main}.
\qed

\begin{remark}[A finite certificate]
Although the digit vectors are modeled as functions
$\N\to\Zmod p$, the proof reads only $2|S|$ carefully selected coordinates.
Thus the construction provides a finite certificate of linear independence.
\end{remark}

\section{Relation to Mathlib's standard digit representation}

The formalization uses the convenient $\Zmod p$-valued definition
\[
  \digit_{p,r}(N)
  =\left[\left\lfloor N/p^r\right\rfloor\right]_p.
\]
Mathlib's standard \texttt{Nat.digits} function returns a finite list of
natural-number digits. After extending that list by zeros and casting its
$r$th entry to $\Zmod p$, the two definitions agree, because canonical casting
to $\Zmod p$ takes the residue modulo $p$. The relevant library statement is
\texttt{Nat.getD\_digits}. A bridge lemma between these presentations is a
natural addition to the public API, allowing later results to be rewritten in
either form without changing the proof construction.

\section{Structure of the Lean formalization}

The formalization is contained in
\href{run:upload/DigitIndependence(1).lean}{\texttt{DigitIndependence(1).lean}}.
Its main components are:

\begin{itemize}[leftmargin=*,itemsep=3pt]
  \item \path{digit} and \path{digits}: one digit and the infinite
        zero-padded digit vector;
  \item \path{pack} and \path{packFin}: block packing for lists and finite
        tuples;
  \item \path{digit_add_high_block} and \path{digit_tail_block}: isolation and
        shifting of base-$p$ blocks;
  \item \path{digit_pack_get} and \path{digit_packFin}: extraction of a digit
        from a selected block;
  \item \path{mul_pack} and \path{mul_packFin}: compatibility of packing with
        multiplication;
  \item \path{exists_crossing}: existence of the least threshold multiplier
        $q_i$;
  \item \path{exists_digits_linearIndependent}: Theorem~\ref{thm:main} for
        every $p\ge2$;
  \item {\small\path{exists_digits_linearIndependent_of_prime}}: compatibility
        corollary for prime bases.
\end{itemize}

In the last step, the Lean proof constructs the linear map $T$ from
Equation~\eqref{eq:T}. The rows of the triangular matrix are the images of the
full digit vectors under $T$. The theorem
\texttt{LinearIndependent.of\_comp} then lifts linear independence from the
images to the original family. To cover composite $p$, the proof applies
\path{Matrix.linearIndependent_cols_of_isUnit} to the transpose of $B$.

\section{Checking the file}

Place the Lean file in a project depending on a compatible recent version of
mathlib and run
\begin{center}
  \texttt{lake env lean 'DigitIndependence(1).lean'}.
\end{center}
Before archival or publication, the exact Lean and mathlib versions should be
pinned in \texttt{lean-toolchain} and \texttt{lakefile.toml}.

\section{Current limitations and possible extensions}

\begin{itemize}[leftmargin=*]
  \item The bounds for $m$, $L$, and therefore $K$ were chosen for a short
        elementary proof, not for efficiency.
  \item The theorem covers every base $p\ge2$, but still uses $\Zmod p$ as the
        coefficient ring. Separating the digit base from the coefficient ring
        would require a different public formulation.
  \item The construction is explicit, but the current theorem states
        existence rather than exposing an efficient executable algorithm.
  \item Natural extensions include sharper bounds, a standard
        \texttt{Nat.digits} corollary, independent base and coefficient ring,
        and a formalized finite-interpolation consequence.
\end{itemize}

\section*{License}

This work is licensed under the
\href{https://creativecommons.org/licenses/by/4.0/}{Creative Commons
Attribution 4.0 International License} (CC BY 4.0). When reusing the work,
please attribute the mathematical idea and construction to Mikhail Katkov,
link to the source and license, and indicate whether changes were made.

\end{document}